\scanpage{58}
\mspage{28}
% archive sheet number: 57

donc si \inserted{$\widetilde Z_0$ est formé d'éléments permutables, p.ex. si} $\operatorname{Comm}_{\unsure{\mathfrak S|\mathscr S}_S}(G_0)$ est commutatif, alors $\widetilde Z_0$ est bien un sous-groupe de $\mathfrak S_S$.

% slanted

\margin{\uline{NB} Si $s\mapsto\unsure{\text{Ét}|\text{Et}}_{\Sigma_0}(s)$ $S\longrightarrow\unsure{\mathscr P|\mathcal P|\mathfrak P}(S)$ est injective, alors pour $\lambda,\mu\in\widetilde Z_0$, on a $\lambda\mu\in\widetilde Z_0\Longleftrightarrow\unsure{\mu\lambda}=\lambda\mu$.}

\uline{Définition} On dit que $A_0$ et $A_1$ sont strictement équivalents, s'ils sont équivalents et isomorphes, i.e. s'il existe \deleted{des} $\lambda,\deleted{\unsure{\in}}f\in\mathfrak S_S$, tels que

\[
 \begin{aligned}
 &\text{a) }\lambda\in\operatorname{Comm}_{\mathfrak S_S}(G_0)\\
 &\text{b) }\forall s,t\in S,\text{ on a}\\
 &\{s,t\}\in A_0\Longleftrightarrow\{f(s),f(\unsure{t|M})\}\in A_1\\
 &\phantom{\{s,t\}\in A_0}\Longleftrightarrow\deleted{\unsure{\{}}\{s,\lambda t\}\in A_1
 \end{aligned}
\]

\deleted{\unsure{La}} C'est encore une relation d'équivalence.

Soit $Z_0$ \deleted{\unsure{$\subset\widetilde Z_0$}} l'ens des \deleted{$\lambda\in\mathfrak S_S$} \inserted{$\lambda\in\widetilde Z_0$} tels que \deleted{il} \deleted{existe} le $A_1$ qu'il définit soit \uline{strictement} équivalent à $A_0$, i.e. soit isom. à $A_0$.

sous réserve que $Z_0$ soit formé d'élément \unsure{permutable|permutables}, $Z_0$ est un sous-groupe de $\mathfrak S_S$. On \unsure{suppose|supposera} dorénavant qu'il en est ainsi ! (i.e. que $Z_0$ est formé d'éléments permutables). \deleted{Alors l'ens}

% slanted

\margin{\uline{NB} Si $s\mapsto\text{Ét}(s)$ $S\longrightarrow\mathscr P(S)$ est injective, alors $\unsure{\mathfrak z|z}=$ ens des éléments involutifs de $\operatorname{Cent}(G_0)$}

Soit

\[
 \mathfrak z_0=Z_0\cap G_0=
 \left\{u\in G_0\ \middle|\ \left\{
 \begin{aligned}
 &\inserted{u\in\operatorname{Cent}(G_0)}\\
 &\deleted{\unsure{\{s,t\}}}\ \text{Ét}_{\Sigma_0}(s)=\text{Ét}_{\Sigma_0}(u^2s)\\
 &\forall s\in S
 \end{aligned}\right.\right\}
 \tag*{(146)}
\]
